\documentclass{article}

\usepackage{polyglossia}
\usepackage{mathtools}
\usepackage[
    math-style=ISO,
    warnings-off={mathtools-colon, mathtools-overbracket},
]{unicode-math}

\setdefaultlanguage{english}

\begin{document}

\section{3.1 Modern Mathematics}

\section{3.2 Single Equation}

Not like this 3-1=0,
but like this \(3-1=2\)
or this 

\[
  3 - 1 = 2
\]

\begin{equation}
    \label{trivial}
    2 + 2 = 4 
\end{equation}
Equation~\ref{trivial}
is true.

\begin{equation}
    \tag{Ingsoc's theorem}
    \label{ingsoc}
    2 + 2 = 5
\end{equation}
Equation~\eqref{ingsoc} is false.

\section{3.2.1 Math Mode}

\(123xyz\) is the same as \(
  1 2 3 x y z
\).\\[1cm]

compare office to \(office\)\\[0.5cm]
\( 1 + \text{one} = 2\)\\[0.5cm]

\( a^2 + b^2 = c^2\)\\[0.5cm]
\(A_x = G_{a\times c}\)

\section{3.3 Building Blocks of Mathematical Formulae}

\section{3.3.1 Basic Arithmetic}

\((4 \times 6) \div 3 = 8\)\\[0.5cm]
\((5 \cdot 3) \divslash 2 \not= 7\)\\[0.5cm]
\(1 \leq 2 \) vs.\
\(1 \leqslant 2 \)\\[0.5cm]

\( a^2 \) is fine, but what about \((a^2)^2\)\\[0.5cm]
\({(a^2)}^2\) is better.\\[0.5cm]

From Homer's theorem it 
follows that
\[
  \sqrt{a} + \sqrt{b} = \sqrt{c}
\]
so we can see that 
\(c = \sqrt{2x^2+7}\)\\[0.5cm]

Find three positive integers 
\(a,b\) and \(c\), such that 
\[
a = \sqrt[7]{b^7 + c^7}
\]\\[0.5cm]

And since 
\[
\frac{a}{b} \leq
\frac{a+b}{b+d} \leq
\frac{c}{d}
\]
it follows that
\(x = \frac{\sqrt{Z+3}}{y^3}\).\\[0.5cm]

\section{3.3.2 Logic and Set Theory}

\[
p \land q \iff 
\lnot (p \implies \lnot q)
\]

\end{document}

